\section[Cote 50 : la formule des traces dans le cadre axiomatique]{Cote 50 : la formule des traces dans le cadre axiomatique\protect\verifie{Folder50}}\label{sec:50}

La cote 50, \emph{Sur bouquin Lang : notes manuscrites}, que l'inventaire date d'après 1965, compte dix pages. Les pages~6 à~10, à lire dans l'ordre 10, 9, 8, 7, puis le haut de 8 et enfin 6, posent pour les endomorphismes des variétés abéliennes un cadre sans points ni diviseurs : une catégorie additive à dualité $D$, un groupe $N(A)$ d'éléments antisymétriques de $\mathrm{Hom}(A, D(A))$, deux cônes $N^+$ et $N^>$, une forme d'intersection $I$ $r$-linéaire à valeurs entières, strictement positive sur $N^>$, et un degré $\nu$ tel que $I(u^*\xi_1, \dots, u^*\xi_r) = \nu(u)\, I(\xi_1, \dots, \xi_r)$ (\lot{50}{1}{7}). La page~8 en tire les coefficients $\sigma_i(\alpha)$ de $\nu(\alpha + n)$ et la formule des traces $\mathrm{tr}(\alpha) = r\, I(\xi^{r-1}, D_\xi(\alpha))/I(\xi^r)$, où $D_\xi(\alpha) = (\alpha + 1)^*\xi - \alpha^*\xi - \xi$ ; la page~6 l'écrit pour $\alpha\beta'$ et $\alpha\alpha'$, puis s'arrête sur le mot \q{Corollaire}, que la lecture complète par la positivité de l'involution de Rosati, $\mathrm{tr}(\alpha\alpha') > 0$ pour $\alpha \neq 0$ (\lot{50}{1}{6}).

Pour une courbe elliptique $E$ et un endomorphisme $\alpha$, le degré de $\alpha + n$ est un polynôme en l'entier $n$, de la forme $n^2 + \mathrm{tr}(\alpha)\,n + \deg\alpha$ ; pour $\alpha = m$ entier, on trouve $(n + m)^2$, d'où $\mathrm{tr}(m) = 2m$. La page cherche à obtenir ce polynôme et sa trace à partir des seuls axiomes, sans géométrie. La trace s'y lit comme un rapport de deux nombres d'intersection. La positivité de Rosati, si elle suivait des axiomes, donnerait la structure de l'algèbre $\mathrm{End}(A) \otimes \mathbf Q$. La vérification ci-dessous montre que la formule des traces suit des axiomes à condition que $I$ soit symétrique, et que la positivité n'en suit pas.

\noindent\emph{Conventions.} On fixe un objet $A$ et une classe $\xi$ dont le morphisme $\varphi_\xi$ est un $\mathbf Q$-isomorphisme, et l'on transporte tout dans $E = \mathrm{End}(A)$ par $\psi(\eta) = \varphi_\xi^{-1}\varphi_\eta$ : l'involution de Rosati $\alpha' = \varphi_\xi^{-1}\,{}^t\alpha\,\varphi_\xi$ est une anti-involution de $E$, et $\varphi_{\alpha^*\eta} = {}^t\alpha\,\varphi_\eta\,\alpha$ se lit $\psi(\alpha^*\eta) = \alpha'\psi(\eta)\alpha$, avec $\psi$ injectif et $\psi(\xi) = 1$. Les cônes, quand ils servent, sont stables par addition et par multiplication par $n > 0$, $N^> \subset N^+$, $N^>$ engendre $N$, $\varphi_x$ est un $\mathbf Q$-isomorphisme pour $x \in N^>$, $I > 0$ sur $(N^>)^r$, $u^*N^+ \subset N^+$ pour tout $u$ et $u^*N^> \subset N^>$ si $\nu(u) \neq 0$.

\begin{enonce}[Pages 8 et 9, le polynôme $\nu(\alpha + n)$]
Les identités (1)--(3) de la page~9 valent : $\psi(D_\xi(\alpha, \beta)) = \alpha'\beta + \beta'\alpha$, $D_\xi(\alpha, \beta) = D_\xi(\alpha'\beta) = D_\xi(\beta'\alpha)$, $D_\xi(\alpha') = D_\xi(\alpha)$ ; de plus $(uv)^* = v^*u^*$ et $\nu(uv) = \nu(u)\nu(v)$. Il existe un unique polynôme de $\mathbf Q[X]$ qui prend en tout entier $n$ la valeur $\nu(\alpha + n)$ ; il est unitaire de degré $2r$, de terme constant $\nu(\alpha)$, et les $\sigma_i(\alpha)$ sont ses coefficients.
\end{enonce}

\noindent L'énoncé tient. Les théorèmes \texttt{nu\_mul}, \texttt{charPoly\_unique} et \texttt{charPoly\_monic} sont dans le fichier \href{https://github.com/Commutator-IO/grothendieck-project/blob/main/lean/Grothendieck/Folder50.lean}{\texttt{Folder50.lean}}.

\begin{encadre}
\begin{theoreme}[D'après la page 8, avec $I$ symétrique\piste{50-axiomatic-trace-formula}]
Si $I$ est symétrique, $\mathrm{tr}(\alpha)\, I(\xi^r) = r\, I(\xi^{r-1}, D_\xi(\alpha))$ pour tout $\alpha \in E$. En particulier $\mathrm{tr}(\alpha'\alpha)\, I(\xi^r) = 2r\, I(\xi^{r-1}, \alpha^*\xi)$ et $\mathrm{tr}(\alpha\alpha')\, I(\xi^r) = 2r\, I(\xi^{r-1}, \alpha'^*\xi)$.
\end{theoreme}
\end{encadre}

\begin{proof}
On a $\psi((\alpha + n)^*\xi) = (\alpha' + n)(\alpha + n) = n^2 + n(\alpha + \alpha') + \alpha'\alpha$, donc $(\alpha + n)^*\xi = n^2\xi + n\,D_\xi(\alpha) + \alpha^*\xi$ par injectivité de $\psi$. Par l'axiome du degré, $\nu(\alpha + n)\, I(\xi^r) = I\bigl((n^2\xi + n\,D_\xi(\alpha) + \alpha^*\xi)^r\bigr)$, que la multilinéarité développe en une somme sur les $f : \{1, \dots, r\} \to \{\xi, D_\xi(\alpha), \alpha^*\xi\}$, de poids $2$, $1$, $0$. Le poids total $2r$ n'est atteint que pour $f$ constante en $\xi$ ; le poids $2r - 1$ l'est pour les $r$ fonctions qui prennent une fois la valeur $D_\xi(\alpha)$, et la symétrie ramène chacun des $r$ termes à $I(\xi^{r-1}, D_\xi(\alpha))$. Enfin $\psi(D_\xi(\alpha'\alpha)) = 2\alpha'\alpha$, d'où $D_\xi(\alpha'\alpha) = 2\alpha^*\xi$.
\end{proof}

\begin{remarque}
La page~6 écrit $\alpha^*\xi$ en regard de $\mathrm{tr}(\alpha\alpha')$ ; les axiomes donnent $\alpha^*\xi$ pour $\alpha'\alpha$ et $\alpha'^*\xi$ pour $\alpha\alpha'$. L'égalité des deux formes est $\mathrm{tr}(\alpha\alpha') = \mathrm{tr}(\alpha'\alpha)$, que le cadre ne fournit pas sans autre argument.
\end{remarque}

\begin{contrexemple}[Contre-exemple de l'édition, $I$ non symétrique]
Soient $E = N = \mathbf Z^2$, d'involution triviale, $\psi = \mathrm{id}$, $\alpha^*\eta = \alpha^2\eta$, $\xi = (1, 1)$, $r = 2$, $I(x, y) = x_1y_2$, $\nu(u) = u_1^2u_2^2$, $N^+ = \mathbf N^2$, $N^> = (\mathbf N \setminus \{0\})^2$. Tous les axiomes sont satisfaits, et pour $\alpha = (1, 0)$ on a $\mathrm{tr}(\alpha)\, I(\xi^2) = 2$, tandis que $2\, I(D_\xi(\alpha), \xi) = 4$ et $2\, I(\xi, D_\xi(\alpha)) = 0$.
\end{contrexemple}

\begin{proof}
$I(u^2x, u^2y) = u_1^2u_2^2\, x_1y_2$, $I$ est $> 0$ sur $(N^>)^2$, $(1, 1)$ et $(1, 2)$ engendrent $\mathbf Z^2$, $x\cdot(x_2, x_1) = x_1x_2(1, 1)$, et $u^*$ préserve les deux cônes. Ici $\nu(\alpha + n) = (1 + n)^2n^2 = n^4 + 2n^3 + n^2$, donc $\mathrm{tr}(\alpha) = 2$, et $D_\xi(\alpha) = (4, 1) - (1, 0) - (1, 1) = (2, 0)$.
\end{proof}

\noindent L'axiome (e) dit seulement $I$ $r$-linéaire. Dans ce modèle, la formule échoue que $D_\xi(\alpha)$ soit placé en premier ou en second.

\begin{contrexemple}[Contre-exemple de l'édition, le Corollaire de la page 6]
Soient $E = \mathbf H[\mathbf Z, 0, 0, 0]$, l'algèbre de quaternions où $i^2 = j^2 = 0$, munie de sa conjugaison, $N = \mathbf Z$, $\psi(x) = x$, $\alpha^*x = (\mathrm{re}\,\alpha)^2x$, $\xi = 1$, $r = 1$, $I(x) = x$, $\nu(\alpha) = (\mathrm{re}\,\alpha)^2$, $N^+ = \mathbf N$, $N^> = \mathbf N \setminus \{0\}$. Tous les axiomes sont satisfaits, ainsi que $I(\xi^{r-1}, \eta) > 0$ pour tout $\eta \in N^+$ non nul, mais $i \neq 0$ et $\mathrm{tr}(i'i) = 0$.
\end{contrexemple}

\begin{proof}
$\alpha'x\alpha = x\,\alpha'\alpha = (\mathrm{re}\,\alpha)^2x$, car $\alpha'\alpha = (\mathrm{re}\,\alpha)^2$ quand $i^2 = j^2 = 0$. Comme $i'i = -i^2 = 0$, $\nu(i'i + n) = n^2$ et $\mathrm{tr}(i'i) = 0$.
\end{proof}

\noindent La lecture déduit le Corollaire de ce que $\alpha^*\xi$ est dans le cône positif ; ici $i^*\xi = 0$. La piste proposait d'ajouter la positivité de $I$ sur $N^+$. Elle tient dans ce modèle, et le Corollaire y est faux.

\begin{enonce}[Énoncé de l'édition, le Corollaire sous deux axiomes de plus]
Si $I$ est symétrique, si $I(\xi^{r-1}, \eta) > 0$ pour tout $\eta \in N^+$ non nul et si $\alpha^*\xi \neq 0$ pour tout $\alpha \neq 0$, alors $\mathrm{tr}(\alpha'\alpha) > 0$ pour tout $\alpha \neq 0$.
\end{enonce}

\noindent Voir \texttt{tr\_star\_mul\_self\_pos} : $\alpha^*\xi \in N^+$ est non nul, et $I(\xi^r) > 0$.

Ne sont pas démontrés ici : la catégorie elle-même et la Propriété fondamentale~1 sur les produits, qui n'est \q{purement formelle} que si $N$ contient tous les éléments antisymétriques ; la première moitié du dossier (pages~2 et~4).
